\documentclass[11pt]{article}
\usepackage[margin=1in]{geometry}
\usepackage[T1]{fontenc}
\usepackage{lmodern}
\usepackage{microtype}
\usepackage{amsmath,amssymb,amsthm}
\usepackage{booktabs,array}
\usepackage{enumitem}
\usepackage[round]{natbib}
\usepackage[hidelinks]{hyperref}
\setlength{\emergencystretch}{2em}

% ---- notation
\newcommand{\R}{\mathbb{R}}
\renewcommand{\P}{\mathbb{P}}
\newcommand{\E}{\mathbb{E}}
\newcommand{\one}{\mathbf{1}}
\newcommand{\N}{\mathrm{N}}
\newcommand{\abs}[1]{\lvert #1\rvert}
\newcommand{\norm}[1]{\lVert #1\rVert}
\newcommand{\ind}[1]{\mathbf{1}\{#1\}}
\newcommand{\Cov}{\operatorname{Cov}}
\newcommand{\Var}{\operatorname{Var}}
\newcommand{\Corr}{\operatorname{Corr}}
\newcommand{\sd}{\operatorname{sd}}
\newcommand{\lmin}{\lambda_{\min}}
\newcommand{\yr}{y^{\mathrm r}}
\newcommand{\Xr}{X^{\mathrm r}}
\newcommand{\bstar}{\beta^*}
\newcommand{\tstar}{\theta^*}
\newcommand{\sigs}{\sigma_*}
\newcommand{\Sz}{S_0}
\newcommand{\sz}{s_0}
\newcommand{\Sig}{\Sigma_n}
\newcommand{\Om}{\Omega}
\newcommand{\pis}[1]{\pi(#1\mid y)}
\newcommand{\PIP}{\mathrm{PIP}}
\newcommand{\HPM}{\mathrm{HPM}}
\newcommand{\MPM}{\mathrm{MPM}}
\newcommand{\Shat}{\hat S}
\newcommand{\Zhat}{\hat Z}
\newcommand{\tSz}{\tilde S_0}
\newcommand{\rn}{\sqrt{\log n/n}}
\newcommand{\OP}{O_{\P}}
\newcommand{\oP}{o_{\P}}
\newcommand{\Dn}{D_n}
\newcommand{\Hl}{\mathcal H}
\newcommand{\Acal}{\mathcal A}
\newcommand{\Ccal}{\mathcal C}

% ---- statements
\theoremstyle{plain}
\newtheorem{thm}{Theorem}
\newtheorem{prop}[thm]{Proposition}
\theoremstyle{definition}
\newtheorem{asm}{Assumption}
\renewcommand{\theasm}{A\arabic{asm}}
\newtheorem{planned}{Planned result}
\renewcommand{\theplanned}{P\arabic{planned}}
\newcommand{\notproved}{\textbf{Planned (not proved).}\ }
\setlist[enumerate]{leftmargin=2em,itemsep=1pt,topsep=3pt}
\setlist[itemize]{leftmargin=1.5em,itemsep=1pt,topsep=3pt}

\begin{document}
\begin{center}
{\LARGE Theory for SMC-BVS:\\[3pt] consistency of the posterior and of the reported model}\\[8pt]
{\large Statements in the format of Section~2 of Song and Liang (2015) \quad$\cdot$\quad 25 November 2026}
\end{center}
\medskip

% =====================================================================================================
\section{Introduction}\label{sec:intro}

SMC-BVS is a method for Bayesian variable selection in the linear model when the number of predictors $p$ is much
larger than the sample size $n$. Its model and priors are standard. Its sampler is new: a sequential Monte Carlo
sampler built to find the posterior mass when predictors are correlated. This document states the theory of the method
in the format of Section~2 of \citet{SongLiang2015}, whose split-and-merge method is the closest to ours in the kind of
theory it proves. It contains statements only. The full statements, with every constant, are in the project's
specification document, which we call the specification.

The document is organised as follows. Section~\ref{sec:model} gives the model, the true density, the asymptotic regime,
the priors and the posterior. Section~\ref{sec:sampler} describes the sampler and what it reports. In
Section~\ref{sec:selection} we show that the posterior concentrates on the true model at the rate $\rn$, so that the
inclusion probabilities, the highest posterior model and the median probability model are all consistent
(Theorem~\ref{thm:post}). In Section~\ref{sec:density} we show that the density fitted by the posterior converges to the
true density in Hellinger distance at the same rate, with no lower bound on the coefficients (Theorem~\ref{thm:dist}).
These two results are about the exact posterior, which cannot be computed at the sizes we target. Section~\ref{sec:software}
is about the software. In Section~\ref{sec:reported} we show that the model a run reports is the true model, whatever
the seed (Theorem~\ref{thm:reported}). In Section~\ref{sec:screening} we show that a small reported model certified as a
local optimum contains every true predictor, under conditions that allow strong common factors
(Theorem~\ref{thm:screen}). In Section~\ref{sec:evidence} we show that the evidence estimate is unbiased
(Proposition~\ref{prop:unbiased}). In Section~\ref{sec:sharp} we show that the rate of Theorem~\ref{thm:post} cannot be
improved (Theorem~\ref{thm:sharp}). Section~\ref{sec:planned} lists five planned results; none of them is proved.

Table~\ref{tab:guide} lists the results and what each one assumes. Assumptions A1--A10 are about the data. Conventions
C1--C3 idealise the computer. Each is stated where it is first needed.

\begin{table}[ht]
\centering\small
\caption{The results, what each is about, and what each assumes.}\label{tab:guide}
\setlength{\tabcolsep}{4pt}
\begin{tabular}{@{}l>{\raggedright\arraybackslash}p{0.25\linewidth}>{\raggedright\arraybackslash}p{0.22\linewidth}>{\raggedright\arraybackslash}p{0.30\linewidth}@{}}
\toprule
Result & About & Design & Assumes \\
\midrule
Theorem~\ref{thm:post} & exact posterior: selection & random, sub-Gaussian rows & A1--A3, A4(a), A5(a), A6 \\
Theorem~\ref{thm:dist} & exact posterior: fitted density & random, i.i.d.\ pairs & A7 \\
Theorem~\ref{thm:reported} & reported model & random, Gaussian rows & A2, A4(b), A5(b), A6, A8, A9; C1 (C2 for the last clause) \\
Theorem~\ref{thm:screen} & certified reported model & fixed & A1, A10; C1 \\
Proposition~\ref{prop:unbiased} & evidence estimate & any data set & C3 \\
Theorem~\ref{thm:sharp}(a) & exact posterior: lower bound & any data set & nothing \\
Theorem~\ref{thm:sharp}(b) & exact posterior: lower bound & orthogonal or Gaussian null predictors & A1, A3, and the conditions in the statement \\
\bottomrule
\end{tabular}
\end{table}

% =====================================================================================================
\section{The model, the regime and the priors}\label{sec:model}

\subsection{The data and the true density}\label{sec:truth}

Let $\Dn$ denote a data set of $n$ observations: a response $\yr\in\R^n$ and an $n\times p$ matrix of predictors
$\Xr$. The superscript means ``raw''. The theorems assume that the data come from a Gaussian linear model with an
intercept,
\begin{equation}\label{eq:model}
  \yr=\alpha^*\one_n+\Xr\bstar+\varepsilon,\qquad \varepsilon\sim\N_n(0,\sigs^2I_n),
\end{equation}
where $\one_n$ is the vector of ones, $\alpha^*\in\R$ is the intercept, $\bstar\in\R^p$ is the vector of regression
coefficients and $\sigs^2$ is the noise variance. The \emph{true model} is $\Sz=\{j:\beta^*_j\neq0\}$, of size $\sz$. Its
members are the \emph{true predictors}; the others are \emph{null predictors}.

\begin{asm}[Gaussian linear model]\label{a:gl}
The data follow~\eqref{eq:model}, with a fixed noise variance $\sigs^2>0$ and noise independent of the predictors.
\end{asm}

Assumption~A1 says that the response is a linear function of a few predictors plus independent Gaussian noise of
constant variance. It is model~(1) of \citet{SongLiang2015} with an intercept added. It is reasonable for a continuous
response with roughly symmetric errors, such as a log expression level. It fails for counts, for heavy-tailed errors,
and when the error variance changes with the mean.

In Theorems~\ref{thm:post}--\ref{thm:reported} the design is random. When the rows of $\Xr$ are independent copies of a
vector $x=(x_1,\dots,x_p)^\top$ with law $\nu_x$, the observations are $n$ independent copies of a pair $(x,y)$, and the
true density of $(x,y)$ is
\begin{equation}\label{eq:fstar}
  f^*(x,y)\,\nu_x(dx)\,dy=\phi\big(y;\ \alpha^*+x^\top\bstar,\ \sigs\big)\,\nu_x(dx)\,dy,
\end{equation}
where $\phi(\cdot\,;\mu,\sigma)$ is the Gaussian density with mean $\mu$ and standard deviation $\sigma$. In
Theorem~\ref{thm:screen} the design is fixed, and probability refers to the noise alone.

The method does not work on the raw data. It centres the response and every predictor and scales each to sample
variance one. We write $y$ and $X$ for the standardised versions: every column of $X$ has mean zero and squared norm
$n-1$, and $y^\top y=n-1$.\footnote{A predictor that is constant in the sample is dropped first, and $p$ then denotes
the number of predictors kept. Theorems~\ref{thm:post} and~\ref{thm:dist} say how they treat dropped predictors.} On the
standardised scale the true coefficients are $\theta^*_j=\beta^*_j\times(\text{sample standard deviation of predictor }j)$,
and $\tstar=(\theta^*_1,\dots,\theta^*_p)^\top$ has the same support $\Sz$.

\subsection{The asymptotic regime and sparsity}\label{sec:regime}

As in \citet{SongLiang2015}, we study the method as $n\to\infty$, with $p=p_n$ allowed to grow with $n$. The design and
the coefficients may also change with $n$. A constant is \emph{fixed} if it does not depend on $n$; every constant
introduced in an assumption is fixed. For random variables, $V_n=\OP(a_n)$ means that $V_n/a_n$ is bounded in
probability, and $V_n=\oP(a_n)$ that it tends to zero in probability. We write $\norm{v}$ for the Euclidean norm of a
vector, $\ind{E}$ for the indicator of an event, and $I$ for an identity matrix.

\begin{asm}[Dimension]\label{a:dim}
$2\le p\le n^A$ for a fixed $A\ge1$.
\end{asm}

Assumption~A2 lets the number of predictors exceed the sample size by any fixed power: $p=10^6$ with $n=200$ is
$A\approx2.6$. It excludes $p$ that grows faster than every power of $n$, such as $p=e^{\sqrt n}$. The reason is a race.
By chance, the best of $p$ null predictors gains about $\log p$ in log-likelihood, and the price that the prior and the
marginal likelihood put on one more predictor, about $\log p+\tfrac12\log n$, must stay ahead of that gain.
\citet{SongLiang2015} allow $p$ to grow almost exponentially, at the price of a slower rate, but they require $p$ to
grow faster than some power of $n$. We allow small $p$ as well.

\begin{asm}[Sparsity]\label{a:sp}
The number of true predictors $\sz\ge1$ is fixed.
\end{asm}

Assumption~A3 says that a fixed, small number of predictors matter, however large $n$ and $p$ become. It is the strict
sparsity condition~(4) of \citet{SongLiang2015}. It holds for a trait driven by a handful of genes. It fails for a trait
with hundreds of small effects, more of which become detectable as $n$ grows. For our rate it is nearly necessary:
by Theorem~\ref{thm:sharp}(a), the rate $\rn$ is possible only if $\sz=O(\sqrt{\log n})$ (for $\sz\le\min(K-1,p/2)$).

\subsection{Prior specification}\label{sec:prior}

A \emph{model} is a set $S\subseteq\{1,\dots,p\}$ of predictors, of size $\abs{S}$, and $X_S$ holds the columns of $X$
in $S$. Given $S$, the method fits
\[
  y\mid S,\theta_S,\sigma^2\ \sim\ \N_n(X_S\theta_S,\ \sigma^2I_n),
\]
with coefficients $\theta_S\in\R^{\abs{S}}$ and variance $\sigma^2>0$. There is no intercept because the data are
centred. The model space $\Om$ is the set of all models with at most $K=\min(p,\lfloor n/2\rfloor)$ predictors. The
priors are
\[
  \pi(S)\ \propto\ \binom{2p-1}{\abs{S}}^{-1}\ (S\in\Om),\quad
  \theta_S\mid S,\sigma^2\ \sim\ \N\big(0,\ n\sigma^2(X_S^\top X_S+I)^{-1}\big),\quad
  \sigma^2\ \sim\ \mathrm{IG}(0.01,\,0.01).
\]
The model prior is the beta-binomial prior restricted to $\Om$: each predictor enters independently with probability
$q$, and $q\sim\mathrm{Beta}(1,p)$. Going from $k$ to $k+1$ predictors costs a factor of about $(k+1)/(2p)$ in prior
odds, so the price of one more predictor grows with $p$. This is the multiplicity control of \citet{ScottBerger2010}.
The coefficient prior is the $g$-prior of \citet{Zellner1986} with $g=n$, with a unit ridge added so that it stays proper
when columns are collinear. The variance prior is a diffuse inverse gamma law.

\citet{SongLiang2015} use a prior of the same three parts: independent inclusion with a fixed probability and a cap on
the model size, a Gaussian coefficient prior, and an inverse gamma variance prior. Their inclusion probability and their
cap are sequences to be specified. Here nothing is left to tune: mixing over $q$ replaces the first, and the
cap is $K$.

\subsection{The posterior and its summaries}\label{sec:posterior}

The coefficients and the variance integrate out in closed form. With $G_S=X_S^\top X_S$, $z_S=X_S^\top y$ and
$H_S=(n+1)G_S+I$, the marginal likelihood of $S$ is
\[
  m(S)\ \propto\ \Big(\frac{\det H_S}{\det(G_S+I)}\Big)^{-1/2}\big(0.02+y^\top y-n\,z_S^\top H_S^{-1}z_S\big)^{-(n/2+0.01)},
\]
and the posterior over models is
\begin{equation}\label{eq:posterior}
  \pis{S}=\frac{\pi(S)\,m(S)}{Z},\qquad Z=\sum_{S\in\Om}\pi(S)\,m(S).
\end{equation}
In $Z$, $\pi(S)$ is normalised over $\Om$ and $m(S)$ includes its constant factor, given in the specification, so $Z$ is the evidence itself. We call $\pi(S)m(S)$ the \emph{score} of $S$ and $Z$ the \emph{evidence}. The usual summaries are the inclusion
probabilities $\PIP_j=\sum_{S\ni j}\pis{S}$, the highest posterior model $\HPM=\arg\max_S\pis{S}$, which
\citet{SongLiang2015} call the maximum a posteriori model, and the median probability model $\MPM=\{j:\PIP_j>1/2\}$ of
\citet{BarbieriBerger2004}.

% =====================================================================================================
\section{The sampler and what it reports}\label{sec:sampler}

\subsection{The sampler in brief}\label{sec:steps}

The posterior~\eqref{eq:posterior} is known up to the constant $Z$. The difficulty is to find where its mass sits in an
enormous model space. Correlated predictors make this hard in two ways. A true predictor and a close proxy give two
models with nearly the same score. And a null predictor that is correlated with a true one looks useful on its own,
because it carries part of the true predictor's signal; we say it \emph{leaks}.

Sequential Monte Carlo starts many models, called \emph{particles}, from an easy distribution close to the posterior
and moves them to the posterior in small steps. Our easy distribution, the \emph{start law}, treats small blocks of
correlated predictors exactly. That settles the choice between a true predictor and its proxies inside a block. What the
start law cannot see is leakage between blocks. A greedy search from the mode of the start law removes leaked
predictors one at a time. The steps are as follows.
\begin{enumerate}
  \item \emph{Blocks.} Cluster the predictors by their sample correlations into blocks $C_1,\dots,C_J$ of at most
    $B=\min(\lfloor\sqrt n\rfloor,\lfloor\sqrt p\rfloor,16)$ predictors. For every subset $U$ of a block $C_k$, compute
    $h_k(U)=m(U)/m(\emptyset)$.
  \item \emph{Start law.} Set $\pi_0(S)=\pi(S)\prod_{k}h_k(S\cap C_k)/\bar Z_0$ for $S\in\Om$, where $\bar Z_0$ is the
    normalising constant. This is the posterior one would have if the marginal likelihood factorised over blocks. A
    recursion over the blocks gives $\bar Z_0$, exact draws from $\pi_0$, and its mode $\tSz$.
  \item \emph{Path to the posterior.} Let $\Delta(S)=\log\{m(S)/m(\emptyset)\}-\sum_k\log h_k(S\cap C_k)$, the part of the
    log marginal likelihood that comes from interactions between blocks. The path is $\pi_\xi\propto\pi_0\,e^{\xi\Delta}$
    for $0\le\xi\le1$, and $\pi_1$ is the posterior. A \emph{calibration} stage, made of a pilot run and a few diagnostic
    runs, chooses a grid $0=\xi_0<\xi_1<\dots<\xi_T=1$ and the number of particles $N$. Both are then frozen.
  \item \emph{Particles.} Four independent \emph{islands} each draw $N$ particles from $\pi_0$. At step $t$ an island
    gives particle $S_v$ the weight $\omega^{(t)}_v=\exp\{(\xi_t-\xi_{t-1})\Delta(S_v)\}$, resamples the particles in
    proportion to these weights, and moves each particle by Markov moves that leave $\pi_{\xi_t}$ invariant: a Gibbs
    update of every predictor, a random swap, swaps with correlated neighbours, and one proposal drawn afresh from
    $\pi_0$.
  \item \emph{Greedy searches.} A search moves from the current model to its best neighbour (add one predictor, delete
    one, or swap one) as long as the score strictly increases, up to a cap on the number of moves. Two searches are run:
    one from the best model the particles found, with cap $\lceil\log_2(n+p+1)\rceil$, and one from the start mode $\tSz$,
    with cap $\lceil\log_2(n+p+1)\rceil+\abs{\tSz}$. The searches use no random numbers.
\end{enumerate}

\subsection{What a run reports}\label{sec:outputs}

A run reports four things.
\begin{enumerate}[label=(\roman*)]
  \item \emph{The reported model} $\Shat$: the model with the highest score among all models the run scored, ties going
    to the smaller model, then to the lexicographically smaller one. It is the run's estimate of the $\HPM$. It is \emph{certified} if it is the end point of one of
    the two searches and that search found no neighbour with a higher score. A certified model is a local optimum of
    the posterior under additions, deletions and swaps.
  \item \emph{Posterior model probabilities, up to $Z$:} the list of at most $40p+1$ best models scored, with their scores.
    By~\eqref{eq:posterior} the ratio of two scores is the exact posterior odds of the two models.
  \item \emph{Inclusion probabilities:} for each predictor $j$, the estimate $\widehat{\PIP}_j$ is the average, over
    particles and islands at $\xi=1$, of the conditional probability that $j$ is in the model given the rest of the
    particle. The estimated median model is $\{j:\widehat{\PIP}_j>1/2\}$.
  \item \emph{The evidence estimate:}
    \begin{equation}\label{eq:zhat}
      \Zhat=\frac14\sum_{r=1}^{4}\Zhat_r,\qquad
      \Zhat_r=m(\emptyset)\,\bar Z_0\prod_{t=1}^{T}\Big\{\frac1N\sum_{v=1}^N\omega^{(t)}_v\Big\}\quad\text{for island }r,
    \end{equation}
    with $\omega^{(t)}_v$ the weights of island $r$ at step $t$.
\end{enumerate}
When $\Om$ is small enough the method enumerates it instead. Everything it reports is then exact, and $\Shat$ counts as
certified. A run can also stop without output, for instance when it would exceed its time or memory allowance. We say
that a run \emph{produces output} when it does not stop.

Theorems~\ref{thm:reported} and~\ref{thm:screen} are about $\Shat$, and Proposition~\ref{prop:unbiased} is about
$\Zhat$. No result in this document is about the estimated inclusion probabilities or the list; the specification has
separate results for them. For $p$ near $10^6$ a parallel version of the sampler is used. The results here are stated
for the standard version, not for that one.

\subsection{The idealised computer}\label{sec:conventions}

The results about the software idealise the computer in three ways. These conventions are not assumptions about the
data. The specification lists them item by item.
\begin{itemize}
  \item[\textbf{C1}] \emph{The method.} Exact arithmetic, the default settings, and a valid input and seed.
  \item[\textbf{C2}] \emph{The machine.} C1, and a machine that never overflows or fails, ideal random numbers, and at
    least the memory and time that the method's own formulas ask for. Two fixed integer limits on the search caps are
    also ignored; no input that the software can represent reaches them.
  \item[\textbf{C3}] \emph{The sampler.} Exact arithmetic; ideal random numbers that, in the final run, are independent
    of the data and of the calibration; exact draws from the start law and an exact $\bar Z_0$; and Markov moves that do
    what their definitions say.
\end{itemize}
C1 fixes what ``the method'' is. C2 is needed only to say that every run finishes, since a real run can be stopped by
its time or memory limit. C3 is needed for Proposition~\ref{prop:unbiased}, which is an exact identity and so needs
exact sampling. Rounding error and imperfect random number generators are outside the results.

% =====================================================================================================
\section{Consistency of the posterior}\label{sec:exact}

\subsection{Selection of the true model}\label{sec:selection}

To select the true model, the posterior needs three things beyond A1--A3: predictors with light tails, true effects
that do not vanish, and no small group of predictors that is nearly collinear.

\begin{asm}[Law of the rows]\label{a:rows}
The rows of $\Xr$ are independent and identically distributed, with mean $m_n$ and covariance matrix $\Sig$, which is a
correlation matrix. In addition, one of:
\begin{enumerate}[label=(\alph*)]
  \item \emph{sub-Gaussian on sparse directions:} for a fixed $L>0$, a row $x$ satisfies
    $\E\exp\{\lambda u^\top(x-m_n)\}\le\exp\{L^2\lambda^2u^\top\Sig u/2\}$ for every $\lambda\in\R$ and every $u\in\R^p$
    with at most $5(A+1)(\sz+1)\log n$ nonzero entries;
  \item \emph{Gaussian:} the rows are $\N_p(m_n,\Sig)$. This is case (a) with $L=1$.
\end{enumerate}
\end{asm}

Assumption~A4 controls the tails of the predictors: every combination of a few predictors has tails no heavier than a
Gaussian's. With it, the sample correlations of all small groups of predictors are close to the population ones at the
same time. Case (a) holds for Gaussian rows, and for rows $x=m_n+\Lambda\xi$ with $\Lambda\Lambda^\top=\Sig$ when the coordinates of
$\xi$ are independent with mean zero, variance one and a common bound. It fails for heavy-tailed predictors and for a
binary predictor whose frequency tends to zero. Condition~1 of \citet{SongLiang2015} plays the same role: it asks for
Gaussian-like tails of averages of up to a prespecified number of predictors. Theorem~\ref{thm:reported} asks for
Gaussian rows, case (b), because A8 and A9 are written with partial covariances, which describe conditional dependence
fully only for Gaussian vectors.

\begin{asm}[Fixed signals]\label{a:fs}
One of:
\begin{enumerate}[label=(\alph*)]
  \item \emph{raw scale:} $\alpha^*$ and $\bstar$ are deterministic, and $b_{\mathrm{lo}}\le\abs{\beta^*_j}\le
    b_{\mathrm{hi}}$ for every $j\in\Sz$, with fixed $0<b_{\mathrm{lo}}\le b_{\mathrm{hi}}<\infty$;
  \item \emph{standardised scale:} $\yr=\alpha\one_n+X\tstar+\varepsilon$, where $X$ is the standardised design,
    $\varepsilon\sim\N_n(0,\sigs^2I_n)$ is independent of $\Xr$, $\sigs^2>0$ is fixed, $\tstar$ is deterministic with a
    fixed number $\sz\ge1$ of nonzero entries, and $\theta_{\mathrm{lo}}\le\abs{\theta^*_j}\le\theta_{\mathrm{hi}}$ for
    every $j\in\Sz$, with fixed $0<\theta_{\mathrm{lo}}\le\theta_{\mathrm{hi}}<\infty$.
\end{enumerate}
\end{asm}

Assumption~A5 says that each true effect is bounded away from zero and from infinity, by constants that do not change
with $n$. A true predictor is selected only if it improves the fit by more than the price of one more predictor, which
is of order $\log p$; a fixed effect improves it by an amount of order $n$. The assumption fails for effects that shrink
with $n$, for example of order $\rn$. Form (a) fixes the coefficients in the original units. Form (b) fixes them on the
scale the method works on; it contains the linear model and the sparsity assumption, so it takes the place of A1 and A3
in Theorem~\ref{thm:reported}. \citet{SongLiang2015} state no lower bound on the nonzero coefficients.

\begin{asm}[Eigenvalue floor]\label{a:efs}
$\lmin(\Sigma_{n,S})\ge\varphi^2$ for every set $S$ of at most $5(A+1)(\sz+1)\log n$ predictors, for a fixed
$\varphi\in(0,1]$. Here $\Sigma_{n,S}$ is the sub-matrix of $\Sig$ for the predictors in $S$, and $\lmin$ is the
smallest eigenvalue.
\end{asm}

Assumption~A6 says that no small group of predictors, of up to a multiple of $\log n$, is close to linearly dependent
in the population. A null predictor that could almost reproduce a true one would give a model with almost the same
score, and the posterior could not separate the two. A floor on the whole matrix, $\lmin(\Sig)\ge\varphi^2$, is enough:
the smallest eigenvalue is $1$ for independent predictors, $1/3$ for autoregressive correlation $0.5$, and $1-\rho$ when
all pairs have correlation $\rho$. There is no upper bound on the eigenvalues, so a strong common factor is allowed. The
assumption fails when a predictor is duplicated, or when the correlation of some pair tends to one. The identifiability
condition of \citet{SongLiang2015} (their Condition~3) plays the same role. For Gaussian predictors it holds when the
smallest eigenvalue of the covariance matrix is of larger order than the square of their rate, a gap that may tend to
zero. Ours is a fixed gap, asked of small sets only.

Under these conditions we have the following theorem.

\begin{thm}[Posterior concentration]\label{thm:post}
Assume A1, A2, A3, A4(a), A5(a) and A6. Then, under the joint law of the design and the noise,
\[
  1-\pis{\Sz}=\OP\big(\rn\big),\qquad
  \max_{1\le j\le p}\big|\PIP_j-\ind{j\in\Sz}\big|=\OP\big(\rn\big),
\]
and $\P(\HPM=\MPM=\Sz)\to1$. Here $p$ in A2, A4(a) and A6 is the number of predictors in the data, before any predictor
that is constant in the sample is dropped, and $\Sz$ is a set of these predictors; the prior, the cap $K$ and the
posterior use the predictors that are kept. A dropped predictor is given $\PIP_j=0$; $\pis{\Sz}$ is taken to be $0$ if
$\Sz$ is not in $\Om$; and if no predictor is kept, the posterior is a point mass at the empty model.
\end{thm}

Theorem~\ref{thm:post} says that the posterior puts all but about $\rn$ of its mass on the true model. Every inclusion
probability is within that distance of $0$ or $1$, and both standard point estimates equal the true model with
probability tending to one. This holds for any correlation structure with an eigenvalue floor, common factors included,
and for any polynomial $p$. Theorem~3 of \citet{SongLiang2015} gives one half of the statement about inclusion
probabilities: that of every identifiable true predictor tends to one, and, as they note, the inclusion probabilities
of null predictors are not controlled. Consistency of their maximum a posteriori model is in their supplement.
Theorem~\ref{thm:post} gives both halves, uniformly over the $p$ predictors and with a rate.

\subsection{Consistency of the fitted density}\label{sec:density}

The first result of \citet{SongLiang2015} is not about selection. It says that the densities drawn from the posterior
are close to the true density $f^*$ in Hellinger distance (their Theorems~1 and~2). We prove a result of the same form.
It needs less than Theorem~\ref{thm:post}, because it asks only that the fitted density be close to the truth, not that
the right predictors be selected. Its setting is stated as one assumption.

\begin{asm}[Setting for the density result]\label{a:dist}
The observations are $n$ independent copies of a pair $(x,y)$ with $x\in\R^p$. Fix an integer $s\ge0$ and constants
$A\ge1$, $L\ge1$, $c_\psi>0$, $\bar b>0$. For a level $M\ge1$, let $q_n=\min\{p,\lceil C_qM\log n\rceil\}$ with
$C_q=2(s(A+2)+4)/\log2$.
\begin{enumerate}[label=(\alph*)]
  \item \emph{Model.} $y=\alpha^*+x^\top\bstar+\varepsilon$ with $\varepsilon\sim\N(0,\sigs^2)$ independent of $x$ and
    $\sigs^2>0$. The quantities $\alpha^*$, $\bstar$, $\sigs$ and the law of $x$ may depend on $n$.
  \item \emph{Sparsity and dimension.} At most $s$ entries of $\bstar$ are nonzero, and $2\le p\le n^A$.
  \item \emph{Sub-Gaussian predictors.} Every predictor has positive variance. With
    $\zeta_j=(x_j-\E x_j)/\sd(x_j)$ and $\zeta=(\zeta_1,\dots,\zeta_p)^\top$,
    $\E\exp(\lambda u^\top\zeta)\le\exp\{L^2\lambda^2\Var(u^\top\zeta)/2\}$ for every $\lambda\in\R$ and every $u\in\R^p$
    with at most $q_n$ nonzero entries.
  \item \emph{No near-exact collinearity.} $\lmin(\Corr(x_U))\ge c_\psi/n$ for every set $U$ of at most $q_n$
    predictors, where $x_U$ is the sub-vector of $x$ for the predictors in $U$.
  \item \emph{Bounded signal.} $\abs{\beta^*_j}\,\sd(x_j)\le\bar b\,\sigs$ for every $j$.
\end{enumerate}
Here $\sd$, $\Var$ and $\Corr$ are population standard deviations, variances and correlation matrices.
\end{asm}

Items (a)--(c) are A1--A4 with less asked: the noise level and the coefficients may change with $n$, and no coefficient
has to be large, or even nonzero. The floor in (d) tends to zero, so it excludes only predictors that are almost exactly
collinear. Item (e) says that no single predictor carries an unbounded signal compared with the noise. A real data set
is inside once exact duplicates and exact linear dependences are removed, for example dummy variables that cover every
level of a factor.

We now define the distance. A posterior draw $(S,\beta,\sigma^2)$, mapped back to the original units, defines the
conditional law $\N(\bar y+(x-\bar x)^\top\beta,\sigma^2)$ of $y$ given $x$, where $\bar x$ and $\bar y$ are the sample
means. The law of $x$ is not estimated. Let $f$ be the joint density of $(x,y)$ when $x$ has its true law and $y$ given
$x$ has this fitted law, and let $f^*$ be the true density~\eqref{eq:fstar}. The Hellinger distance is
\[
  d_H(f,f^*)=\Big(1-\int\!\!\int\sqrt{f(x,y)\,f^*(x,y)}\ \nu_x(dx)\,dy\Big)^{1/2}.
\]
It is the distance of \citet{SongLiang2015} divided by $\sqrt2$. Let $\P^*$ denote the law of the data set $\Dn$ and
$\Pi(\cdot\mid\Dn)$ the posterior of $(S,\beta,\sigma^2)$.

\begin{thm}[Consistency of the fitted density]\label{thm:dist}
There are constants $c_H>0$, $M_0$ and $n_0$, depending only on $(s,A,L,c_\psi,\bar b)$, such that the following holds.
Fix $M\ge M_0$ and assume A7 with this $M$. Then for every $n\ge n_0$ with $\sqrt n/\log n\ge M$,
\[
  \P^*\Big[\ \Pi\Big(d_H(f,f^*)>\sqrt{M\log n/n}\ \Big|\ \Dn\Big)>n^{-c_HM}\ \Big]\ \le\ n^{-c_HM}.
\]
Here $p$ in A7 is the number of predictors in the data, and the method runs on those that are not constant in the
sample. If the posterior is not defined, because fewer than two predictors are left or the response is constant in the
sample, the inner probability is read as $1$.
\end{thm}

Theorem~\ref{thm:dist} says that the density the posterior fits to $y$ given $x$ is within $\sqrt{M\log n/n}$ of the
truth, except on a set of posterior mass $n^{-c_HM}$ and a set of data sets of the same probability. It has the form
of Theorems~1 and~2 of \citet{SongLiang2015} with their rate $\epsilon_n$ equal to $\sqrt{M\log n/n}$: their bound
$\exp(-c_1n\epsilon_n^2)$ is then $n^{-c_1M}$. The rate is explicit, and it holds for the prior of
Section~\ref{sec:prior}, with nothing chosen by the user. The theorem needs no lower bound on the coefficients: effects
too weak to be selected do no harm to the fitted density. Theorem~\ref{thm:post} is about selection; this one is about
prediction.

% =====================================================================================================
\section{Consistency of what the software reports}\label{sec:software}

Theorems~\ref{thm:post} and~\ref{thm:dist} describe a posterior that nobody can compute exactly at this size. The
results of this section are about the output of a run. Section~2 of \citet{SongLiang2015} has no result of this kind:
its theorems are about the exact posterior.

\subsection{The reported model}\label{sec:reported}

The start law sees interactions only inside blocks of at most $16$ predictors. The next two assumptions are therefore
about how predictors relate once up to $15$ others are held fixed. Let $W=(W_1,\dots,W_p)^\top\sim\N_p(0,\Sig)$ be a
centred copy of one row of the design, and let $W_U$ be its sub-vector for a set $U$ of predictors. Then
$\Cov(W_j,W_l\mid W_U)$ is the partial covariance of predictors $j$ and $l$ given those in $U$, and
$\bar\beta_{U,l}=\Sigma_{n,UU}^{-1}\Sigma_{n,Ul}$ is the population regression coefficient of $W_l$ on $W_U$, where
$\Sigma_{n,UU}$ and $\Sigma_{n,Ul}$ are the corresponding blocks of $\Sig$.

\begin{asm}[Visibility]\label{a:vis}
For every true predictor $j\in\Sz$ and every set $U\not\ni j$ of at most $15$ predictors,
$\abs{\Cov(W_j,\,W^\top\tstar\mid W_U)}\ge a$, for a fixed $a>0$.
\end{asm}

Assumption~A8 says that each true predictor stays related to the signal $W^\top\tstar$, by a fixed margin, whichever
$15$ other predictors are held fixed. The start law looks at one block at a time, so it can pick up a true predictor
only through what that predictor explains given the other members of its block. The assumption holds for independent
predictors, for positive autoregressive correlation with true effects of one sign, and for block designs with the true
predictors in different blocks. It fails when other predictors cancel a true predictor's signal: two true predictors
with coefficients $1$ and $-1$ and correlation $\rho$ each have covariance $1-\rho$ with the signal, which vanishes as
$\rho\to1$. The marginal relevance condition of \citet{SongLiang2015} (their Condition~5) asks that every true
predictor have nonzero covariance with the response on its own. Ours asks for a margin, and after conditioning.

\begin{asm}[Small leakage set]\label{a:leak}
Fix $b\in(0,\tfrac12)$ and $b_H>0$ with $15b_H+b<\tfrac12$.
\begin{enumerate}[label=(\alph*)]
  \item \emph{Leakage set.} There is a deterministic set $Q\supseteq\Sz$ with $\abs{Q}\le n^b$ such that:
    \begin{enumerate}[label=(\roman*),leftmargin=1.8em]
      \item \emph{no leak outside $Q$:} $\abs{\Cov(W_l,W_j\mid W_U)}\le\varpi_n$ for every $l\notin Q$, every set $U$ of
        at most $15$ predictors and every $j\in\Sz\setminus U$, where $\varpi_n$ is a deterministic sequence with
        $\varpi_n\sqrt{n/\log n}\to0$;
      \item \emph{locality:} every $l\notin Q$ has a deterministic neighbourhood $\Hl_l\supseteq\{l\}\cup\Sz$ with
        $\abs{\Hl_l}\le n^{b_H}$ such that, for every $U\not\ni l$ with $\abs{U}\le15$ and $U_0=U\cap\Hl_l$,
        $\abs{\Cov(W_l,W_v\mid W_{U_0})}\le\delta_n$ and $\abs{\Cov(W_j,W_v\mid W_{U_0})}\le\delta_n$ for every
        $v\in U\setminus\Hl_l$ and $j\in\Sz\setminus U_0$, where $\delta_n$ is a deterministic sequence with
        $\delta_n\to0$;
      \item \emph{room under the size cap:} for large $n$, $\abs{Q}\le K$, and either $Q$ is the set of all predictors or
        $K\ge2\abs{Q}+32\sz+1$.
    \end{enumerate}
  \item \emph{Floor.} $\lmin(\Sigma_{n,Q\cup\{l\}})\ge\varphi^2$ for every predictor $l$, with $\varphi$ as in A6.
  \item \emph{Local regressions.} There are deterministic sequences $d_n=o(\log n)$ and $L_n\to0$ with
    $L_n^2\log(\abs{Q}+1)=o(\log n)$, and for every predictor $l$ a deterministic set $N_l\subseteq Q\setminus\{l\}$ of at
    most $d_n$ predictors, such that $\norm{\bar\beta_{S,l}-\bar\beta_{S\cap N_l,l}}_1\le L_n$ for every model $S$ with
    $\Sz\subseteq S\subseteq Q$ and $l\notin S$. (The second coefficient vector is filled with zeros to the length of the
    first, and $\norm{\cdot}_1$ is the sum of absolute values.)
\end{enumerate}
\end{asm}

Assumption~A9 limits leakage. Part (a) says that only a small set $Q$ of predictors, fewer than $\sqrt n$, can imitate
part of the signal: every predictor outside $Q$ is nearly uncorrelated with the true ones given any $15$ others, and
conditioning on far-away predictors changes nothing that matters. Part (b) says that no predictor is close to a linear
combination of the members of $Q$. Part (c) says that the regression of any predictor on a model between $\Sz$ and $Q$
is close to its regression on a few members of $Q$. For independent predictors $Q=\Sz$. For block designs $Q$ is the
union of the blocks that contain true predictors, and for autoregressive designs it is a neighbourhood of size
$C\log n$ around them. With a pervasive common factor every null predictor is correlated with every true one, so $Q$
would have to be everything, and A9 fails once $p$ is larger than $\sqrt n$. \citet{SongLiang2015} have no counterpart;
they remark that their method relies on only a small proportion of the predictors being correlated with the true ones.

Under these conditions the reported model is the true model.

\begin{thm}[The reported model equals the true model; simplified, full version in the specification, Theorem~3]\label{thm:reported}
Let $\Sig$ be positive definite and assume A2, A4(b), A5(b), A6, A8, A9 and C1. Then, under the joint law of the design
and the noise, there are events $\mathcal G_n$, determined by the data alone, with $\P(\mathcal G_n)\to1$, on which
every run that produces output reports $\Shat=\Sz$, whatever its seed and whatever its calibration chose. In particular
\[
  \P\big(\text{the run produces output and }\Shat\neq\Sz\big)\to0 .
\]
Under C2 every run produces output, for every $n\ge3$ that the method accepts, and so $\P(\Shat=\Sz)\to1$.
\end{thm}

Theorem~\ref{thm:reported} says that, for almost every data set, the model the software prints is exactly the true
model, and that this does not depend on the seed. The proof follows the design of the sampler. By visibility, the mode
of the start law contains every true predictor. Anything else in it is a leaked predictor, and so belongs to $Q$. The
greedy search from that mode deletes the leaked predictors one at a time and stops at $\Sz$, and no other model the run
scores beats $\Sz$. The search is deterministic, which is why the seed does not matter. The result covers independent,
block and autoregressive designs. It does not cover pervasive common factors, because A9 fails there. The statement is
simplified in one respect: the specification states C2 in full and treats the two integer limits separately.

\subsection{Certified screening}\label{sec:screening}

The next result asks much less of the design and concludes less. It needs neither visibility nor a leakage set. Its
assumption is on the observed design, which is taken as fixed.

\begin{asm}[Conditions for certified screening]\label{a:scr}
The design is fixed. Let $\hat R_S=X_S^\top X_S/(n-1)$ be the sample correlation matrix of the predictors in $S$. For a
\emph{size budget} $k_n\in\{0,\dots,K-1\}$:
\begin{enumerate}[label=(\alph*)]
  \item $\lmin(\hat R_S)\ge\kappa$ for every $S$ with $\abs{S}\le k_n+\sz$;
  \item $\min_{j\in\Sz}\abs{\theta^*_j}\ge\underline\theta>0$ and $\norm{\tstar}_1\le M_1$;
  \item $(k_n+1)\log p+\log n\le c\,n$, with $c=\kappa^2\underline\theta^2/\{200(M_1^2+\sigs^2)\}$.
\end{enumerate}
The constants $\kappa>0$, $\underline\theta$, $M_1$ and $\sigs^2$ are fixed.
\end{asm}

Assumption~A10 is an eigenvalue floor and signal bounds of the same kind as A6 and A5, now on the observed design and
on sets of the budget size. Item (c) lets the budget be as large as a small multiple of $n/\log p$. The true model size
$\sz$ need not be fixed, and common factors are allowed. The assumption fails if some $k_n+\sz$ observed predictors are
nearly collinear, or if true effects are tiny.

\begin{thm}[Certified screening]\label{thm:screen}
Let the design be fixed and assume A1, A10 with size budget $k_n$, and C1. Assumption A3 is not needed. Then
\[
  \P\big(\text{the run produces output},\ \Shat\text{ is certified},\ \abs{\Shat}\le k_n,\ \text{and }\Sz\not\subseteq\Shat\big)\to0 .
\]
\end{thm}

Theorem~\ref{thm:screen} says that a local optimum of the posterior with at most $k_n$ predictors cannot miss a true
predictor, except with vanishing probability. If it missed one, bringing that predictor in would raise the score by
more than chance can explain. So a small certified model can be read as a screening set: it contains every true
predictor. This is a sure screening property, like Theorem~3 of \citet{SongLiang2015}, but for the model a run reports
and not for exact inclusion probabilities. It covers designs with strong common factors, where
Theorem~\ref{thm:reported} does not apply. The theorem bounds the probability of one joint event. It does not say that a
run will certify a small model, it is not a statement conditional on certification, and a certified model may contain
null predictors.

\subsection{The evidence estimate}\label{sec:evidence}

The last result about the software concerns the evidence estimate $\Zhat$ of~\eqref{eq:zhat}. It holds for every data
set, with or without a true model. Let $\Acal$ be the event that the method uses the sampler and its calibration
succeeds, and let $\Ccal$ be the record of everything the calibration chose. For each island, we consider the estimate
$\Zhat_r$ that the island would produce if it ran to the end without interruption.

\begin{prop}[The evidence estimate is unbiased]\label{prop:unbiased}
Assume C3. For every fixed data set, almost surely on $\Acal$,
\[
  \E\big[\Zhat_r\mid\Ccal\big]=Z\quad(r=1,\dots,4),\qquad\text{and hence}\qquad \E\big[\Zhat\mid\Ccal\big]=Z ,
\]
where the expectation is over the random numbers of the final run. Under the joint law of the data and the method's
random numbers, the same identities hold conditionally on the data and $\Ccal$, and
$\E[\ind{\Acal}\Zhat]=\E[\ind{\Acal}Z]$. No condition on the design is needed.
\end{prop}

Proposition~\ref{prop:unbiased} says that the grid and the number of particles can be tuned on the data without
biasing the estimate. They are frozen before the final run draws a single random number, so the classical unbiasedness
of the sequential Monte Carlo estimate of a normalising constant \citep{DelMoral2004} survives the tuning. The
statement is about $\Zhat$, not $\log\Zhat$.

% =====================================================================================================
\section{The rate cannot be improved}\label{sec:sharp}

Theorem~\ref{thm:post} gives the rate $\rn$ for the posterior mass outside the true model. The next theorem shows that
this rate is a property of the posterior, not of the proof.

\begin{thm}[The rate cannot be improved]\label{thm:sharp}
\textup{(a)} \emph{Every design.} For every $n\ge3$, every standardised design with $p\ge2$, every response $y$ and
every model $S$ with $\abs{S}\le\min(K-1,p/2)$,
\[
  \pis{S}\ \le\ \Big(1+\frac{\abs{S}+1}{3\sqrt n}\Big)^{-1}.
\]

\textup{(b)} \emph{Designs that attain the rate.} Assume A1 and A3. Let the true columns of the raw design and $\tstar$
be deterministic, with $\norm{X\tstar}^2\le nM_\mu$ and $\norm{\tstar}^2\le nM_\theta$ for fixed $M_\mu$, $M_\theta$. (If
they are random, the conditions are asked conditionally on them, on events of probability tending to one, with the
same constants.) Assume one of: \textup{(i)} the null columns of $X$ are orthogonal to each other and to the true
columns; or \textup{(ii)} the raw null columns have independent $\N(0,1)$ entries, independent of the raw true columns
and of the noise, and $p\le n^A$ for a fixed $A$. Then, for every $p>\sz$, bounded or growing,
\[
  \E\big[1-\pis{\Sz}\big]\ \ge\ (1-o(1))\,(\sz+1)\,\frac{2(p-\sz)}{2p-1-\sz}\,
  \sqrt{\frac{\log\!\big(p\sqrt n/(\sz+1)\big)}{\pi n}}\qquad(\pi=3.14\ldots),
\]
and, if $p\to\infty$, $1-\pis{\Sz}\ge(1-\oP(1))\,(\sz+1)\sqrt{\log p/(\pi n)}$.
\end{thm}

Part (a) holds for every data set and has no assumption. Under this prior no model of that size, the true one included,
can have posterior mass closer to one than about $(\abs{S}+1)/(3\sqrt n)$. So, for $\sz\le\min(K-1,p/2)$, no design gives a rate faster than
$n^{-1/2}$, and the rate $\rn$ of Theorem~\ref{thm:post} is possible only when $\sz=O(\sqrt{\log n})$. This is why A3 is
nearly necessary. Part (b) shows that in the most favourable designs, with orthogonal or independent null predictors,
the posterior still misses the true model by a constant times $\sqrt{\log p/n}$. With $p=n^A$ this is the rate of
Theorem~\ref{thm:post}, logarithm included. On other designs only the order $n^{-1/2}$ of part (a) is proved.

% =====================================================================================================
\section{Planned results}\label{sec:planned}

None of the following is proved. Each is a target, and the choices still open are named.

\begin{planned}[Standard samplers get trapped where ours does not]
\notproved There is a family of designs whose posterior has several well-separated modes, on which (i) every Markov
chain that adds or removes one predictor at a time, and the weighted tempered Gibbs sampler of
\citet{ZanellaRoberts2019}, has a mixing time that grows quickly with the size of the problem, while (ii) one of
Theorems~\ref{thm:post}, \ref{thm:reported} or~\ref{thm:screen} holds on the same designs. Open: the family of designs,
which of the three theorems, the exact class of chains, and the growth rate.

The target must be positioned against two results on fast mixing. \citet{YangWainwrightJordan2016} show that a
Metropolis--Hastings chain with single flips and swaps mixes in a number of steps polynomial in $n$, $p$ and $\sz$,
under eigenvalue and signal conditions. \citet{ZhouYangVatsRobertsRosenthal2022} show
that an informed chain mixes fast when the posterior has a single mode under additions, deletions and swaps, and give
an example where mixing is slow. The family of designs must therefore violate the conditions of both papers while our
consistency result holds on it. No existing result covers the sampler of \citet{ZanellaRoberts2019}.
\end{planned}

\begin{planned}[The theory at $p=10^7$]
\notproved Take the autoregressive design of the planned $p=10^7$ run at $n=200$, continued as a sequence with
$p_n=\lceil n^{A}\rceil$ and $A=\log(10^7)/\log200\approx3.04$. Show that it satisfies the assumptions of our theorems
with this $A$; for Theorem~\ref{thm:reported} this means exhibiting a leakage set $Q$ that satisfies A9. Open: which
theorems are covered, the set $Q$, how the sequence is continued, and whether the target should instead be the list of
results that a single fit with one seed can support. The run itself uses the parallel version of the sampler, which no
theorem covers.
\end{planned}

\begin{planned}[Theorem~\ref{thm:reported} on the raw scale]
\notproved Theorem~\ref{thm:reported} takes the coefficients fixed on the standardised scale, and
Theorem~\ref{thm:dist} takes them fixed on the raw scale; with $\sz\ge1$ no design satisfies both. The target is to
restate A5(b), A8 and A9 for fixed raw coefficients and to show, under A7 with $s\ge1$, that with probability tending to
one every completed run returns a list of models $\mathcal T$ with posterior mass
$\pi(\mathcal T\mid\Dn)\ge n^{-c_HM/2}$. The Hellinger bound of Theorem~\ref{thm:dist} would then hold for the posterior
restricted to the list the software returns. Open: the raw-scale form of the assumptions, and whether to aim for the
full conclusion $\Shat=\Sz$ on the raw scale.
\end{planned}

\begin{planned}[Theorem~\ref{thm:reported} under an aggregated leak]
\notproved The conclusions of Theorem~\ref{thm:reported} hold when the bound A9(a)(i), which is on each predictor
outside $Q$ separately, is replaced by a bound on the leaks of all predictors outside $Q$ taken together. Open: the
form of the aggregated condition.
\end{planned}

\begin{planned}[A wrong model, with $p$ fixed]
\notproved Drop the linear model altogether. Let the rows $(x_i,y_i)$ be independent draws from any fixed law on
$\R^p\times\R$ with $p$ fixed and $\Sigma=\Var(x)$ positive definite. Let $\beta^0=\Sigma^{-1}\Cov(x,y)$ and
$\alpha^0=\E y-\E x^\top\beta^0$ be the best linear predictor, $r=y-\alpha^0-x^\top\beta^0$ its residual, $v=\E
r^2>0$, and $S^0=\{j:\beta^0_j\neq0\}$. Let $f_0$ be the joint density with $x$ from its true law and $y\mid
x\sim\N(\alpha^0+x^\top\beta^0,v)$, and let $f$ be the density fitted by a posterior draw, as in
Theorem~\ref{thm:dist}. Assume $\E[\norm{x}^2r^2]<\infty$ and $\E r^4<\infty$. Target:
\[
  1-\pi(S^0\mid\Dn)=\OP(n^{-1/2}),\qquad \max_j\big|\PIP_j-\ind{j\in S^0}\big|=\OP(n^{-1/2}),
\]
$\P(\HPM=\MPM=S^0)\to1$, and, for every $M_n\to\infty$,
\[
  \Pi\Big(d_H(f,f_0)>\frac{M_n}{\sqrt n}\ \text{ or }\ \mathrm{KL}(f_0\,\|\,f)>\frac{M_n}{n}\ \Big|\ \Dn\Big)\to0\quad\text{in probability,}
\]
where $\mathrm{KL}(f_0\,\|\,f)=\int f_0\log(f_0/f)$ is the Kullback--Leibler divergence. Here $f_0$ is the
Kullback--Leibler projection of the true law onto the Gaussian linear family, whenever that divergence is finite for
some member of the family. So without Gaussian errors, a linear conditional mean or constant error variance, the
posterior would select the support of the best linear predictor and concentrate on the closest Gaussian linear model.

This is the counterpart of Theorem~4 of \citet{SongLiang2015}, their result for misspecified models. They keep the
linear model as the truth, fit a fixed set of candidate predictors that misses some true ones, and show that the
posterior concentrates at the Kullback--Leibler projection $f_0$, with no rate; their Corollary~2 adds that the
selected model contains the support of the projection. The target here allows any truth with the stated moments, and
asks for the exact support and a rate. It is for fixed $p$ only.
\end{planned}

\bibliographystyle{abbrvnat}
\bibliography{../../shared/refs}

\end{document}
